\documentclass[11pt]{article}
\usepackage[english]{babel}
\usepackage[letterpaper,top=1.85cm,bottom=1.9cm,left=2.3cm,right=2.3cm]{geometry}
\usepackage{amsmath,booktabs,graphicx,float}
\usepackage[labelfont=bf,labelsep=period]{caption}
\usepackage[sort,round,semicolon,authoryear]{natbib}
\usepackage[none]{hyphenat}
\usepackage{microtype,xurl}
\usepackage[colorlinks=true,allcolors=blue]{hyperref}
\hypersetup{pdftitle={The Accessibility Cost of Protecting Journeys in Coordinated Bus-Service Cuts},pdfauthor={Adhar Sharma}}
\urlstyle{same}
\emergencystretch=3em
\setlength{\parindent}{0pt}
\setlength{\parskip}{4pt}
\setlength{\bibsep}{2pt}
\renewcommand{\bibfont}{\small}
\raggedbottom
\clubpenalty=10000
\widowpenalty=10000
\renewcommand{\bibsection}{\par\bigskip\noindent\textbf{References}\par\medskip}

\title{The Accessibility Cost of Protecting Journeys in Coordinated Bus-Service Cuts}
\author{Adhar Sharma\\[2pt]\small The University of Tokyo, Japan}
\date{}
\begin{document}
\maketitle
\begin{abstract}
When bus operators cut service, letting one operator take over another's route can improve access while removing journeys retained by the best fixed-operator timetable. We study this choice using 832 departure-time queries in a reconstructed Sapporo network. At a 20\% cut, reassignment raises one-hour reachability from 411 to 415: nine queries are gained and five lost. Preserving all 411 forfeits that net gain of four, whereas at a 45-minute limit the operators can preserve all 326 reference queries and add two. Equally good 415-query plans place the five losses in different communities.
\end{abstract}
\textit{Keywords:} bus-service cuts; accessibility; hold-harmless clause; operator coordination; timetable optimization.

\section{QUESTIONS}
We compare \emph{two network-wide optimized timetables}: one retains each route's current operator, while the other permits Bankei and Jotetsu to exchange operating responsibility for selected routes. Both choose retained departures under separate operator limits; this is not a comparison with independently planned cuts. Sapporo's transport plan calls for operator cooperation amid service contraction \citep{sapporoPlan2024}, and downsizing and capacity-sharing research motivates the experiment \citep{guillot2022,pantelidis2024}.

A planner evaluating a revised, post-cut timetable might wish to promise that every origin--destination--departure-time opportunity in a selected fixed-operator plan remains reachable. We test this \emph{hypothetical} hold-harmless guarantee, not an existing Sapporo contract. How much extra access can route reassignment provide, and how much survives if none of the selected reference journeys may be lost? We also vary the allowed losses ($L$) and permitted reassignments ($K$).

\section{METHODS}
Public GTFS snapshots provide three Bankei route units (48 trips), five Jotetsu units (205 trips), and fixed subway service (678 trips) \citep{bankeiGTFS,jotetsuGTFS,subwayGTFS}. Bus service is dated 9 November 2026, subway service 3 February 2020. Reassignment keeps the same departure times, stops, terminals and turnaround assumptions; no depot or staffing costs are added.

We sample 11 route-end/interchange origins, five subway/interchange destinations and 16 hourly desired departures from 06:00 to 21:00. Three origin--destination pairs have identical endpoints; removing their 48 same-origin-and-destination \emph{queries} leaves 832. For example, can a traveller leaving Nakanosawa at 10:00 reach Makomanai Station by 11:00? The 60-minute limit includes initial waiting and transfers. Queries have equal weight, rather than observed passenger counts; the destinations and hourly starts sample access across the selected corridors \citep{geurs2004,farber2014}.

The model chooses which nested departure packages to retain and, when permitted, who operates each route, maximizing reachable queries. Separate revenue-hour and experimental vehicle caps apply to each operator (four Bankei, 23 Jotetsu vehicles). Within-route vehicle chaining is permitted, but not cross-route chaining. Seven units can change operator. A zero-loss rule preserves the \emph{identities} of all reachable queries in one selected \emph{post-cut} fixed-operator optimum, not all full-service journeys or their original travel times. For a given reference, ``gain'' is the difference between unrestricted and fixed-operator optimal coverage; ``cost'' is unrestricted minus zero-loss coverage. Package construction, exact inputs, selected reference IDs and code are in the supplement and \url{https://github.com/Todai-ops/sapporo-bus} \citep{sharma2026sapporoRepo}.

\section{FINDINGS}
Full service makes 426 of the 832 queries reachable within an hour. With a 20\% cut, a rule-based timetable serves 400, reoptimizing departures with operators fixed serves 411, and allowing reassignment serves 415 (Table~\ref{tab:primary}). Thus better fixed-operator timetabling adds 11 over the rule-based plan; changing operators adds a further \emph{four}. Protecting the particular 411-query optimum forfeits four queries, about 1\% of that reference's coverage and 100\% of the reassignment gain.

\begin{table}[H]
\centering
\caption{Primary 60-minute results, offset zero (reachable queries of 832). ``Keep reference'' preserves each query reachable under the selected post-cut fixed-operator optimum.}
\label{tab:primary}
\small
\resizebox{\linewidth}{!}{%
\begin{tabular}{@{}lrrrrrr@{}}
\toprule
Hours cut & Rule-based & Routes fixed & Routes may change & Keep reference & Extra gain & First useful $L$\\
\midrule
10\% & 407 & 419 & 421 & 419 & 2 & 3\\
20\% & 400 & 411 & 415 & 411 & 4 & 3\\
30\% & 393 & 403 & 406 & 403 & 3 & 3\\
\bottomrule
\end{tabular}}
\end{table}

In the selected 411-query plan, Bankei runs no Komaoka departures. Allowing three losses \emph{first permits a coverage improvement}, not the first transfer: Jotetsu takes 10 Komaoka departures while reducing N96 from 24 to 18. Five queries become reachable and three are lost, yielding 413 within both operators' caps. One loss is the Nakanosawa 10:00--Makomanai journey. With five losses allowed, one reassignment reaches 415 (nine gained, five lost). Two tied 415-query plans both transfer Komaoka: one cuts N96 and loses five Nakanosawa queries, while another cuts the Fujino route and loses five Fujino queries. Their equal totals conceal different choices for communities.

\begin{table}[H]
\centering
\caption{Original offset and fixed-operator reference sensitivity (60-minute threshold). All-gain forfeiture is counted only where reassignment improves coverage.}
\label{tab:phase}
\small
\resizebox{\linewidth}{!}{%
\begin{tabular}{@{}lrrrrr@{}}
\toprule
Cut & Offsets with gain & Offsets forfeiting all & $L^*$ range & Distinct fixed references & Cost share across references\\
\midrule
10\% & 6/8 & 4/6 & 1--3 & 4 & 100\%\\
20\% & 7/8 & 7/7 & 1--3 & 2 & 100\%\\
30\% & 8/8 & 7/8 & 1--8 & 5 & 66.7--100\%\\
\bottomrule
\end{tabular}}
\end{table}

At the same 20\% cut, a \emph{45-minute} criterion selects a different 326-query reference and permits two more queries with none lost; the one-hour criterion permits no extra queries under full protection (Table~\ref{tab:newtests}). At 75 and 90 minutes, reassignment offers no additional coverage. The temporal-spacing package menu has three equally good 414-query references: protecting them yields 414, 415 or 416, against 418 unrestricted. The archived reference is the middle case (75\% cost); across the three, the cost is 50--100\%.

\begin{table}[H]
\centering
\caption{20\%-cut sensitivity. Temporal spacing reports the \emph{archived} reference (protected 415; cost share 75\%); across three optimal references protection yields 414--416 (cost 50--100\%). Cost share is undefined with no reassignment gain.}
\label{tab:newtests}
\small
\resizebox{\linewidth}{!}{%
\begin{tabular}{@{}lrrrrr@{}}
\toprule
Time limit/package choice & Routes fixed & Routes may change & Keep reference & Extra gain & Cost share\\
\midrule
45 min, 8 levels & 326 & 328 & 328 & 2 & 0\%\\
60 min, 8 levels & 411 & 415 & 411 & 4 & 100\%\\
75 min, 8 levels & 529 & 529 & 529 & 0 & n/a\\
90 min, 8 levels & 612 & 612 & 612 & 0 & n/a\\
60 min, 16 levels & 411 & 415 & 412 & 4 & 75\%\\
60 min, temporal spacing & 414 & 418 & 415 & 4 & 75\%\\
\bottomrule
\end{tabular}}
\end{table}

The value of reassignment depends on the journey definition, retained departures and \emph{which} optimal reference is promised. The mixed-year feeds, stylized vehicle caps and unweighted query sample limit application. A practical agreement should identify the actual reference timetable and disclose where access is gained or lost, not just compare aggregate totals.

\bibliographystyle{plainnat}
\bibliography{references}
\end{document}
